% Generated by scripts/assemble.py; corrections belong in source/corrections.json.
\examhead{19}{中国人民大学 2024 年数学分析试卷}{24}
% p025-b01; source PDF page 25
\par\Needspace{5\baselineskip}
\blockanchor{p025-b01}
\majorquestion{一}(15 分) 已知 $n$ 为正整数, 计算极限 $\lim_{n\to \infty}\sum_{i = 1}^{n}\frac{\sin\frac{i\pi}{n}}{n + \frac{i}{n}}.$\par
% p025-b02; source PDF page 25
\par\Needspace{5\baselineskip}
\blockanchor{p025-b02}
\majorquestion{二}(15 分) 设 $n$ 为正整数, 函数 $f(x)$ 在 $[0, +\infty)$ 上一致连续, 满足对任意 $x > 0$, 都有 $\lim_{n \to \infty} f(x + n) = 0$, 证明: $\lim_{x \to +\infty} f(x) = 0$.\par
% p025-b03; source PDF page 25
\par\Needspace{5\baselineskip}
\par\Needspace{6\baselineskip}
\blockanchor{p025-b03}
\majorquestion{三}(15 分) 设函数 $f(x)$ 在 $[0,2]$ 上存在二阶导数, 且 $|f(x)| \leq 1$, $|f''(x)| \leq 1$, $\forall x \in [0,2]$. 证明:\par
\nopagebreak[4]
% p025-b04; source PDF page 25
\blockanchor{p025-b04}
\displayformula{\left| f ^ {\prime} (x) \right| \leq 2, \forall x \in [ 0, 2 ].}
% p025-b05; source PDF page 25
\par\Needspace{5\baselineskip}
\par\Needspace{6\baselineskip}
\blockanchor{p025-b05}
\majorquestion{四}(15 分) 设函数 $f(x) = \int_{1}^{x} \mathrm{e}^{t^2} \, \mathrm{d}t$，证明：存在 $\xi \in (1,2), \eta \in (1,2)$，满足\par
\nopagebreak[4]
% p025-b06; source PDF page 25
\blockanchor{p025-b06}
\displayformula{f (\xi) \doteq (2 - \xi) \mathrm{e} ^ {\xi^ {2}}, f (2) = \eta \mathrm{e} ^ {\eta^ {2}} \ln 2.}
% p025-b07; source PDF page 25
\par\Needspace{5\baselineskip}
\blockanchor{p025-b07}
\majorquestion{五}(15 分) 求函数 $f(x, y, z) = xyz$ 在条件 $x^2 + y^2 + z^2 = 1$, $x + y + z = 0$ 下的极值.\par
% p025-b08; source PDF page 25
\par\Needspace{5\baselineskip}
\par\Needspace{6\baselineskip}
\blockanchor{p025-b08}
\majorquestion{六}(15 分) 设函数 $f(x), g(x)$ 在区间 $[a, b]$ 连续, 且 $f(x)$ 严格递增, $0 \leq g(x) \leq 1$. 证明:\par
\nopagebreak[4]
% p025-b09; source PDF page 25
\blockanchor{p025-b09}
\displayformula{\int_ {a} ^ {a + \int_ {a} ^ {b} g (t) \mathrm{d} t} f (x) \mathrm{d} x \leq \int_ {a} ^ {b} f (x) g (x) \mathrm{d} x.}
% p025-b10; source PDF page 25
\par\Needspace{5\baselineskip}
\par\Needspace{6\baselineskip}
\blockanchor{p025-b10}
\majorquestion{七}(15 分) 设函数 $f(x, y)$ 有连续偏导数, 对任意光滑曲线 $L$, 曲线积分 $\int_{L} 2xy \, \mathrm{d}x + f(x, y) \, \mathrm{d}y$ 与路径无关. 求 $f(x, y)$ 满足对任意实数 $t$ 都有\par
\nopagebreak[4]
% p025-b11; source PDF page 25
\blockanchor{p025-b11}
\displayformula{\int_ {(0, 0)} ^ {(t, 1)} 2 x y \mathrm{d} x + f (x, y) \mathrm{d} y = \int_ {(0, 0)} ^ {(1, t)} 2 x y \mathrm{d} x + f (x, y) \mathrm{d} y.}
% p025-b12; source PDF page 25
\par\Needspace{5\baselineskip}
\blockanchor{p025-b12}
\majorquestion{八}(15 分) 已知 $V$ 是由曲面 $y = x^2 (x \geq 0)$ 和平面 $y = x, x + y + z = 2$ 及 $z = 0$ 围成的区域. 计算积分 $\iiint_{V} x \mathrm{e}^{y + z} \mathrm{d}x \mathrm{d}y \mathrm{d}z$.\par
% p025-b13; source PDF page 25
\par\Needspace{5\baselineskip}
\blockanchor{p025-b13}
\majorquestion{九}(15 分) 证明函数项级数 $\sum_{n=1}^{\infty} \frac{x^n}{1 - x^{2n}} (1 - x) \sin nx$ 在 $\left(\frac{1}{2}, 1\right)$ 上一致收敛.\par
% p025-b14; source PDF page 25
\par\Needspace{5\baselineskip}
\par\Needspace{6\baselineskip}
\blockanchor{p025-b14}
\majorquestion{十}(15 分) 给定 $n$ 为正整数, 设函数\par
\nopagebreak[4]
% p025-b15; source PDF page 25
\blockanchor{p025-b15}
\displayformula{T _ {n} (x) = \frac {\alpha_ {0}}{2} + \sum_ {k = 1} ^ {n} (\alpha_ {k} \cos k x + \beta_ {k} \sin k x).}
% p025-b16; source PDF page 25
\par\Needspace{6\baselineskip}
\blockanchor{p025-b16}
其中 $\alpha_0, \alpha_k, \beta_k$ 为常数, $1 \leq k \leq n$. 证明:\par
\nopagebreak[4]
% p025-b17; source PDF page 25
\blockanchor{p025-b17}
\displayformula{\max _ {- \pi \leq x \leq \pi} \left| T _ {n} ^ {\prime} (x) \right| \leq n ^ {2} \max _ {- \pi \leq x \leq \pi} \left| T _ {n} (x) \right|.}
